% Originally: content/week-04.tex, \section{Taylor expansions} (Corsin Week 4, Friday)

% Source: Corsin Nick/Class Notes/Week 4.pdf, p. 9
\section{Taylor expansions}
\label{sec:taylor_expansions}

% Sonnet 5 (Medium)
\subsection{The general formula}

In the lecture, you have seen the following formula. For
$f \in C^k(U \subseteq \mathbb{R}^n, \mathbb{R})$, $\bar x \in U$:
\[f(\bar x + h) = \sum_{|\alpha| = 0}^{k} \partial^\alpha f(\bar x)\,\frac{h^\alpha}{\alpha!} + o(|h|^k)\]
where $\alpha$ is a \newterm{multi-index} \germanfor{Multiindex} in $\mathbb{N}^n$.

In practice, we never use this formula directly, but it is worth unpacking once (e.g., as in \cref{ex:taylor_maximal_expansion}).

% Generator: Claude Opus 5 (High) --- the multi-index notation above is only well defined
% because mixed partials commute; the fact is used twice below (in the expansion and in
% \cref{rem:second_order_taylor_hessian}) but stated nowhere in the source notes.
% Neither Corsin's notes nor the working formula below mention Schwarz's theorem, yet both depend
% on it. Reclassified from ainote to plain prose: this is a mathematical prerequisite.
% Transition: Claude Opus 5 (High)
Before unpacking the formula, one hypothesis hidden in the notation deserves to be made explicit.
Writing $\partial^\alpha$ for a multi-index $\alpha$ presumes that it does not matter \emph{in
which order} the derivatives are taken, since otherwise $\partial^{(1,1)}$ would be ambiguous
between $\partial_1\partial_2$ and $\partial_2\partial_1$. That the presumption is legitimate is
itself a theorem, and it is stated next.

\begin{theorem}[Schwarz / Clairaut --- mixed partial derivatives commute]
\label{thm:schwarz_mixed_partials}
Let $U \subseteq \mathbb{R}^n$ be open and $f \in C^2(U,\mathbb{R})$. Then for all
$i,j \in \{1,\dots,n\}$ and all $x \in U$,
\[\partial_i\partial_j f(x) = \partial_j\partial_i f(x).\]
More generally, for $f \in C^k(U,\mathbb{R})$ any $k$ partial derivatives may be permuted
freely, so $\partial^\alpha f$ depends only on the multi-index $\alpha$ and not on the order of
differentiation. In particular the Hessian $\Hess f(x)$ of
\cref{rem:second_order_taylor_hessian} is a \textbf{symmetric} matrix.
\end{theorem}

\begin{remark}[The hypothesis is not decoration]
\label{rem:schwarz_needs_c2}
Continuity of the second derivatives is genuinely needed. The standard counterexample is
\[f(x,y) := \begin{cases} xy\,\dfrac{x^2-y^2}{x^2+y^2} & (x,y)\neq(0,0), \\[1ex]
  0 & (x,y)=(0,0),\end{cases}\]
for which a direct computation from the definition gives
$\partial_1\partial_2 f(0,0) = 1$ but $\partial_2\partial_1 f(0,0) = -1$. This $f$ has both
mixed partials everywhere, but they are not continuous at the origin, so
\cref{thm:schwarz_mixed_partials} does not apply, and indeed fails. Since everything in
\cref{sec:taylor_expansions} assumes $f \in C^k$, the issue never arises here. It is worth
knowing that it \emph{could}.
\end{remark}

% Supplement: Sascha Brack/Class Notes/Week_05_Notes_Monday.pdf, p. 7
\begin{exercise}[Applying Schwarz's theorem]
\label{ex:schwarz_application}
Let $f : \mathbb{R}^2 \to \mathbb{R}$ be defined by
\[f(x,y) := x^2 e^{y^2} + e^{-x^2} \sin\left(e^{x^2} \left(x^5 + 3x^2 + 100x + \frac{1}{2+\sin(x)}\right)\right).\]
\begin{enumerate}[label=\textbf{(\alph*)}]
  \item Show that $f \in C^2(\mathbb{R}^2, \mathbb{R})$.
  \item Compute $\partial_2 \partial_1 f(x,y)$. (Make it as easy as possible for yourself.)
\end{enumerate}
\end{exercise}

\begin{exercisesolution}[\cref{ex:schwarz_application}]
\begin{enumerate}[label=\textbf{(\alph*)}]
  \item Every term in $f$ is a composition of $C^2$ functions (polynomials, exponentials, sine),
    and the denominator $2+\sin(x)$ is bounded away from zero. Thus $f$ is $C^2$.
  \item Since $f \in C^2$, by Schwarz's theorem (\cref{thm:schwarz_mixed_partials}) we have
    $\partial_2 \partial_1 f = \partial_1 \partial_2 f$. Computing $\partial_1$ of the second
    term would be tedious. However, notice that the second term depends only on $x$, so its
    derivative with respect to $y$ is $0$. Thus, we first compute $\partial_2 f$:
    \[\partial_2 f(x,y) = x^2 (2y e^{y^2}) + 0 = 2x^2 y e^{y^2}.\]
    Then we compute $\partial_1$ of this result:
    \[\partial_1 \partial_2 f(x,y) = \partial_1 (2x^2 y e^{y^2}) = 4xy e^{y^2}.\]
    By Schwarz's theorem, $\partial_2 \partial_1 f(x,y) = 4xy e^{y^2}$.
\end{enumerate}
\end{exercisesolution}

% Source: Corsin Nick/Class Notes/Week 4.pdf, pp. 9--10
\begin{exercise}[Maximal Taylor expansion in $\mathbb{R}^2$]
\label{ex:taylor_maximal_expansion}
Let $f \in C^3(\mathbb{R}^2,\mathbb{R})$. Write out the maximal Taylor approximation.
\end{exercise}

% Kept inline: solved directly in Corsin Nick/Class Notes/Week 4.pdf, pp. 9--10.
\begin{exercisesolution}[\cref{ex:taylor_maximal_expansion}]
The multi-indices are:
\begin{itemize}
  \item $|\alpha| = 0$: $\alpha = (0,0)$
  \item $|\beta| = 1$: $\beta_1 = (1,0)$, $\beta_2 = (0,1)$
  \item $|\gamma| = 2$: $\gamma_1 = (1,1)$, $\gamma_2 = (2,0)$, $\gamma_3 = (0,2)$
  \item $|\delta| = 3$: $\delta_1 = (3,0)$, $\delta_2 = (2,1)$, $\delta_3 = (1,2)$, $\delta_4 = (0,3)$
\end{itemize}
Recall that
\[\partial^\alpha = \partial_1^{\alpha_1}\partial_2^{\alpha_2}\cdots\partial_n^{\alpha_n}, \qquad \text{e.g.\ } \partial^{\delta_2} = \partial^{(2,1)} = \partial_1^2\partial_2,\]
and
\[\alpha! = \alpha_1!\cdots\alpha_n!, \qquad \text{e.g.\ } \gamma_2! = (2,0)! = 2\cdot 1 = 2.\]
So:
\[
\begin{aligned}
f(\bar x + h) = &\ f(\bar x) && (\alpha) \\
&+ \partial_1 f(\bar x)h_1 + \partial_2 f(\bar x)h_2 && (\beta_1, \beta_2)\\
&+ \partial_1\partial_2 f(\bar x)h_1h_2 + \partial_1^2 f(\bar x)\frac{h_1^2}{2} + \partial_2^2 f(\bar x)\frac{h_2^2}{2} && (\gamma_1, \gamma_2, \gamma_3)\\
&+ \partial_1^3 f(\bar x)\frac{h_1^3}{6} + \partial_1^2\partial_2 f(\bar x)\frac{h_1^2h_2}{2} && (\delta_1, \delta_2)\\
&+ \partial_1\partial_2^2 f(\bar x)\frac{h_1h_2^2}{2} + \partial_2^3 f(\bar x)\frac{h_2^3}{6} && (\delta_3, \delta_4)\\
&+ o(|h|^3)
\end{aligned}
\]
\end{exercisesolution}

\begin{ainote}
Corsin writes the recall line as
$\partial^\alpha = \partial_{\alpha_1}\partial_{\alpha_2}\cdots\partial_{\alpha_n}$ with the
example $\partial^{(2,1)} = \partial_2\partial_1$. Read literally that is the wrong object, since
a multi-index $\alpha$ records \textit{how many times} each $\partial_i$ is applied, so
$\partial^{(2,1)} = \partial_1^2\partial_2$. His own expansion above uses
$\partial_1^2\partial_2 f(\bar x)\,\tfrac{h_1^2h_2}{2}$ for $\delta_2$, confirming the intended
meaning. Corrected above.
\end{ainote}

% Sonnet 5 (Medium)
\subsection{The working formula}

% Source: Corsin Nick/Class Notes/Week 4.pdf, p. 11
We can rewrite this as
\[f(\bar x + h) = \sum_{\ell=0}^{k}\frac{1}{\ell!}\big(h_1\partial_1 + h_2\partial_2\big)^\ell f(\bar x).\]

\begin{theorem}[Taylor working formula]
\label{thm:taylor_working_formula}
For $f \in C^k(U \subseteq \mathbb{R}^n, \mathbb{R})$ with $\bar x + th \in U$ for all
$t \in [0,1]$:
\[f(\bar x + h) = \sum_{\ell=0}^{k}\frac{1}{\ell!}\left(\sum_{m=1}^{n} h_m\partial_m\right)^{\!\ell} f(\bar x) + o(|h|^k)\]
\end{theorem}

\begin{ainote}
Both sums in \cref{thm:taylor_working_formula} are written with lower limit $\ell = 1$ (resp.\
starting index) in the source. The $\ell = 0$ term is $f(\bar x)$ itself, which the explicit
expansion above does include, so the lower limit must be $0$. Corrected above; compare the
$|\alpha| = 0$ line of \cref{ex:taylor_maximal_expansion}, which is the same term written in
multi-index form.
\end{ainote}

% Generator: Claude Sonnet 5 (High) --- new content: \cref{thm:taylor_working_formula} only
% gives the asymptotic (Peano) remainder $o(|h|^k)$. The exact (Lagrange) remainder is genuinely
% different information --- an equality valid for that specific $h$, not just eventually small
% --- and was not stated anywhere in this document.
\begin{theorem}[Taylor's theorem with Lagrange remainder]
\label{thm:taylor_lagrange_remainder}
Let $U \subseteq \mathbb{R}^n$ be open, $f \in C^{k+1}(U,\mathbb{R})$, $\bar x \in U$, and
$h \in \mathbb{R}^n$ with $\bar x+th \in U$ for all $t \in [0,1]$. Then there exists
$\theta \in (0,1)$ such that
\[f(\bar x+h) = \sum_{\ell=0}^{k}\frac{1}{\ell!}\Big(\sum_{m=1}^n h_m\partial_m\Big)^{\!\ell}
  f(\bar x)
  \ +\ \frac{1}{(k+1)!}\Big(\sum_{m=1}^n h_m\partial_m\Big)^{\!k+1}f(\bar x+\theta h).\]
Equivalently, in multi-index notation,
\[f(\bar x+h) = \sum_{|\alpha|\leq k}\partial^\alpha f(\bar x)\,\frac{h^\alpha}{\alpha!}
  \ +\ \sum_{|\alpha|=k+1}\partial^\alpha f(\bar x+\theta h)\,\frac{h^\alpha}{\alpha!}.\]
\end{theorem}

\begin{proof}
Let $\varphi(t) := f(\bar x+th)$ for $t \in [0,1]$, well defined and $C^{k+1}$ since $f \in
C^{k+1}$ and the segment $\bar x+th$ stays in $U$. By induction on $\ell$ via the chain rule,
which is exactly the computation underlying \cref{thm:taylor_working_formula} itself,
\[\varphi^{(\ell)}(t) = \Big(\sum_{m=1}^n h_m\partial_m\Big)^{\!\ell}f(\bar x+th)
  \qquad \text{for } 0 \leq \ell \leq k+1.\]
Apply the one-variable Taylor theorem with Lagrange remainder (\emph{Analysis I}) to $\varphi$
on $[0,1]$: there is $\theta \in (0,1)$ with
\[\varphi(1) = \sum_{\ell=0}^k\frac{\varphi^{(\ell)}(0)}{\ell!} + \frac{\varphi^{(k+1)}(\theta)}{(k+1)!}.\]
Substituting $\varphi(1) = f(\bar x+h)$, $\varphi^{(\ell)}(0) = \big(\sum_mh_m\partial_m\big)^\ell
f(\bar x)$, and $\varphi^{(k+1)}(\theta) = \big(\sum_mh_m\partial_m\big)^{k+1}f(\bar x+\theta h)$
gives the working-formula version. The multi-index version follows by expanding
$\big(\sum_mh_m\partial_m\big)^\ell$ via the multinomial theorem, exactly as in
\cref{ex:taylor_maximal_expansion}.
\end{proof}

\begin{remark}
Unlike \cref{thm:taylor_working_formula}'s $o(|h|^k)$ remainder, this is an \emph{exact}
equality for the specific $h$ in question, not merely a statement about what happens as
$h \to 0$. That is genuinely different information, and it is what is needed whenever a hard,
uniform error bound is wanted rather than an asymptotic one (e.g.\ bounding
$|f(\bar x+h) - T_k(h)|$ over an entire region at once, given a bound on the $(k{+}1)$-th
derivatives there). The price is that $\theta$ depends on $\bar x$, $h$ \emph{and} $k$ together,
and is almost never computable in practice. That is exactly why the rest of this chapter uses the
asymptotic form for actually computing expansions, and reserves this one for error estimates.
\end{remark}

% Quelle: Linus Lüchinger/Slides/Slides-03-16.pdf, slide 12;
%         Linus Lüchinger/Slides/Slides-03-19.pdf, slides 4--6
% Generator: Claude Opus 5 (High)
\begin{corollary}[A polynomial that fits to order $k$ \emph{is} the Taylor polynomial]
\label{cor:taylor_expansion_unique}
Let $U \subseteq \mathbb{R}^n$ be open, $x_0 \in U$, $f \in C^{k+1}(U)$, and let $P$ be a
polynomial of degree at most $k$ with
\[\lvert f(x_0+h) - P(h)\rvert = o\big(\lvert h\rvert^{k}\big) \qquad \text{as } h \to 0.\]
Then $\partial^\alpha f(x_0) = \partial^\alpha P(0)$ for every multi-index with
$\lvert\alpha\rvert \leq k$. Equivalently, $P$ is \emph{the} Taylor polynomial of $f$ of degree
$k$ at $x_0$.
\end{corollary}

\begin{proof}
By \cref{thm:taylor_working_formula}, the Taylor polynomial $T$ of degree $k$ also satisfies
$\lvert f(x_0+h)-T(h)\rvert = o(\lvert h\rvert^k)$. Subtracting, the polynomial $Q := P-T$ has
degree at most $k$ and satisfies $\lvert Q(h)\rvert = o(\lvert h\rvert^k)$. Suppose
$Q \not\equiv 0$ and let $d \leq k$ be the degree of its lowest non-vanishing homogeneous part
$Q_d$. Choosing $v$ with $Q_d(v) \neq 0$ and restricting to the ray $h := tv$,
\[Q(tv) = t^dQ_d(v) + O(t^{d+1}), \qquad\text{so}\qquad
  \frac{\lvert Q(tv)\rvert}{t^{k}} \sim \lvert Q_d(v)\rvert\, t^{d-k} \not\to 0\]
as $t \downarrow 0$, since $d \leq k$. That contradicts $\lvert Q(h)\rvert = o(\lvert h\rvert^k)$,
so $Q \equiv 0$ and $P = T$.
\end{proof}

\begin{remark}[Reading derivatives off an asymptotic expansion]
\label{rem:reading_derivatives_off_expansions}
This is more useful than it looks, because it runs \emph{backwards}. Normally one computes
derivatives in order to produce an expansion; the corollary says that an expansion, however it
was obtained, hands back the derivatives. Given
\[f(x_0+h) = c + a\transp h + \tfrac12 h\transp Bh + o(\lvert h\rvert^2)\]
one may read off $f(x_0) = c$, $\nabla f(x_0) = a$ and $\Hess f(x_0) = B$ without differentiating
anything. This is the same move by which \cref{def:differentiable} recognises a differential from
an $o(\lvert h\rvert)$ estimate, now carried to higher order. It is the standard route to a Hessian
for functions built out of known series.

Two cautions on the bookkeeping, both of which the exercise below turns on. The estimate must be
$o(\lvert h\rvert^{k})$, not $O(\lvert h\rvert^{k})$: an $O(\lvert h\rvert^{k})$ error can
contain a genuine degree-$k$ term and so corrupts the order-$k$ coefficients, leaving only the
derivatives of order $< k$ determined. And the expansion determines derivatives only up to
order $k$, and nothing beyond it.
\end{remark}

% Supplement: Linus Lüchinger/Slides/Slides-03-19.pdf, slide 6
\begin{exercise}[What does an expansion actually tell you?]
\label{ex:reading_off_expansions}
For each of the following, state precisely which values of $f$ and of its derivatives at the
indicated point are determined by the given information, and which are not.
\begin{enumerate}[label=\textbf{(\alph*)}]
  \item $f : \mathbb{R}^3 \to \mathbb{R}$ with
    $f(x_1,x_2,x_3) - x_1x_2 \in O\big(\lvert x\rvert\big)$, at $x = 0$.
  \item $f : \mathbb{R}^2 \to \mathbb{R}$ with
    $f(1+x_1,\,-1+x_2) - 2 - x_1 \in o\big(\lvert x\rvert\big)$, at $(1,-1)$.
  \item $f : \mathbb{R}^2 \to \mathbb{R}$ with
    $f(-1+x_1,\,x_2) + 2 - 2x_1x_2 + 2x_2^2 - 3x_2^3 \in o\big(\lvert x\rvert^3\big)$,
    at $(-1,0)$.
\end{enumerate}
Assume in each case that $f$ is as smooth as the relevant statement requires.
\end{exercise}

% Supplement: Toby Lane/class-document.pdf, §3.4.4
\begin{remark}[Two ways to weight the same expansion]
\Cref{ex:taylor_maximal_expansion} above and \cref{thm:taylor_working_formula} appear to disagree
on how heavily mixed partial derivatives are weighted. Take the coefficient of $h_1^2h_2$, so
$\alpha = (2,1)$ and $\ell = |\alpha| = 3$. The explicit multi-index expansion gives $\tfrac12$
directly. The working formula instead offers $\tfrac{1}{\ell!}(h_1\partial_1+h_2\partial_2)^\ell$,
and expanding that by the multinomial theorem produces
\[
  \frac{\ell!}{\alpha!}\cdot\frac{1}{\ell!} \;=\; \frac{1}{\alpha!}
  \;=\; \frac{1}{2!\,1!} \;=\; \frac12
\]
for the term $h^\alpha\partial^\alpha f$. The two agree, but they arrive by different routes, and
seeing why is worth the paragraph.

The working formula's $\tfrac{1}{\ell!}$ weights \emph{all} $\ell$-fold products of partials
globally, including repeated ones such as $\partial_1\partial_2$, which appears in more than one
order in the multinomial expansion of $(h_1\partial_1+h_2\partial_2)^\ell$. By Schwarz's lemma
(equality of mixed partial derivatives, valid here because $f \in C^k$) all those repeats collapse
onto the same $\partial^\alpha$, and there are exactly $\ell!/\alpha!$ of them. That collapsing is
what turns the global $1/\ell!$ into the per-term $1/\alpha!$ seen in
\Cref{ex:taylor_maximal_expansion}.

In short, $1/\ell!$ weights \emph{paths} through the mixed partial derivatives, whereas
$1/\alpha!$ weights the \emph{distinct} derivatives themselves, once those paths have been
identified with one another.
\end{remark}

% Sonnet 5 (Medium)
\subsection{A worked example}

% Source: Corsin Nick/Class Notes/Week 4.pdf, pp. 11--12
\begin{exercise}[Taylor expansion of $(x+y^2)e^{-x^2-y^2}$]
\label{ex:taylor_explicit_example}
Compute the Taylor expansion of
\[f : \mathbb{R}^2\to\mathbb{R}, \qquad (x,y) \mapsto (x+y^2)e^{-x^2-y^2}\]
around $(0,0)$ up to third order.
\end{exercise}

% Kept inline: solved directly in Corsin Nick/Class Notes/Week 4.pdf, pp. 11--12.
\begin{exercisesolution}[\cref{ex:taylor_explicit_example}, cumbersome solution]
Calculate all partial derivatives up to third order and use the definition. (Please never do
this.)
\end{exercisesolution}

\begin{exercisesolution}[\cref{ex:taylor_explicit_example}, the easy solution]
The substitution and multiplication rules we used in Analysis 1 still apply:
\[
\begin{aligned}
(x+y^2)e^{-x^2-y^2} &= (x+y^2)\left(1 - x^2 - y^2 + o(|x^2+y^2|)\right) \\
&= x + y^2 - x^3 - xy^2 + o\!\left(|(x,y)|^3\right)
\end{aligned}
\]
\end{exercisesolution}

% Supplement: Sascha Brack/Class Notes/Week_05_Notes_Monday.pdf, pp. 9--11
\begin{example}[More Taylor expansions via Analysis 1]
\label{ex:more_taylor_via_analysis1}
The substitution and multiplication rules from Analysis 1 often make expansions much faster than using the formula directly, provided the expansion is around $(0,0)$.
\begin{enumerate}[label=\textbf{(\alph*)}]
  \item Expand $f(x,y) := \cos(x)e^y$ up to second order around $(0,0)$:
  \[f(x,y) = \left(1 - \frac{x^2}{2} + o(|(x,y)|^2)\right)\left(1 + y + \frac{y^2}{2} + o(|(x,y)|^2)\right) = 1 + y - \frac{1}{2}x^2 + \frac{1}{2}y^2 + o(|(x,y)|^2).\]
  \item Expand $f(x,y) := \cos(x+y^2)$ up to second order around $(0,0)$. Setting $z := x+y^2$:
  \[\cos(z) = 1 - \frac{z^2}{2} + o(|z|^2) \implies \cos(x+y^2) = 1 - \frac{(x+y^2)^2}{2} + o(|(x,y)|^2) = 1 - \frac{1}{2}x^2 + o(|(x,y)|^2),\]
  where terms like $xy^2$ and $y^4$ are absorbed into $o(|(x,y)|^2)$ since their total degree is strictly greater than $2$.
  \item Expand $f(x,y) := \sqrt{y} e^{\sqrt{x}}$ to first order around $(1,1)$. Here the formula is better since the expansion point is not the origin.
  We have $f(1,1) = e$.
  $\partial_1 f(x,y) = \sqrt{y} e^{\sqrt{x}} \frac{1}{2\sqrt{x}} \implies \partial_1 f(1,1) = \frac{e}{2}$.
  $\partial_2 f(x,y) = \frac{1}{2\sqrt{y}} e^{\sqrt{x}} \implies \partial_2 f(1,1) = \frac{e}{2}$.
  Writing $h = (h_1,h_2) = (x-1, y-1)$:
  \[f(1+h_1, 1+h_2) = f(1,1) + \partial_1 f(1,1)h_1 + \partial_2 f(1,1)h_2 + o(|h|) = e + \frac{e}{2}h_1 + \frac{e}{2}h_2 + o(|h|).\]
\end{enumerate}
\end{example}

% Sonnet 5 (Medium)
\subsection{Second-order form and the Hessian}

% Source: Corsin Nick/Class Notes/Week 4.pdf, p. 12
\begin{remark}[Second-order Taylor formula with Hessian]
\label{rem:second_order_taylor_hessian}
The second-order Taylor approximation of $f \in C^2(\mathbb{R}^n,\mathbb{R})$ around
$x_0 \in \mathbb{R}^n$ can be expressed compactly as
\[f(x_0+x) = f(x_0) + \nabla f(x_0)\cdot x + \tfrac{1}{2}x\transp \cdot \Hess f(x_0)\cdot x + o(|x|^2)\]
where
\[\Hess f(x_0) = \big(\partial_i\partial_j f(x_0)\big)_{i,j=1,\dots,n}\]
is the \newterm{Hessian matrix} \germanfor{Hesse-Matrix} of $f$.
\end{remark}

% Generator: Claude Opus 5 (High)
This is the form the rest of the course actually uses, and it is worth reading term by term as
a sequence of ever-better approximations to $f$ near $x_0$:
\begin{itemize}
  \item the \textbf{constant} term $f(x_0)$, the crudest guess, exact only at $x_0$;
  \item the \textbf{linear} term $\nabla f(x_0)\cdot x$, which is the tangent plane and hence
    exactly the differential of \cref{def:differentiable};
  \item the \textbf{quadratic} term $\tfrac12 x\transp\Hess f(x_0)\,x$, the first term that sees
    curvature, and the one that decides in \cref{sec:hessian_test} whether a critical point is a
    minimum, a maximum, or a saddle.
\end{itemize}
By \cref{thm:schwarz_mixed_partials} the Hessian is symmetric, which is what makes that
last decision a question about the \emph{signs of its eigenvalues} rather than anything harder.

\begin{center}
% Generator: Claude Opus 5 (High) --- FIG-W04-05: the three approximations, in one variable,
% where they can actually be drawn. f(x) = 1 + 0.55x - 0.5x^2 + 0.28x^3 about x_0 = 0, so
% f(0) = 1, f'(0) = 0.55, f''(0) = -1: constant 1, tangent 1 + 0.55x, parabola 1 + 0.55x - 0.5x^2.
\begin{tikzpicture}[>=Latex, scale=1.25]
  \draw[->] (-1.7,0) -- (2.1,0) node[right, font=\footnotesize] {$x$};
  \draw[->] (0,-0.3) -- (0,2.6) node[above, font=\footnotesize] {$y$};
  \draw (0,0.05) -- (0,-0.05) node[below, font=\scriptsize] {$x_0$};

  % order 0
  \draw[gray!70, thick, dashed] (-1.6,1) -- (2.0,1);
  \node[gray!70!black, font=\scriptsize, anchor=west] at (2.02,1.0) {order $0$};

  % order 1: tangent
  \draw[green!50!black, thick, domain=-1.6:2.0, samples=2]
    plot (\x, {1 + 0.55*\x});
  \node[green!50!black, font=\scriptsize, anchor=west] at (2.02,2.1) {order $1$};

  % order 2: parabola
  \draw[orange!90!black, thick, domain=-1.3:1.85, samples=60, smooth]
    plot (\x, {1 + 0.55*\x - 0.5*\x*\x});
  \node[orange!90!black, font=\scriptsize, anchor=west] at (1.88,0.40) {order $2$};

  % the function itself
  \draw[blue!75!black, very thick, domain=-1.35:1.75, samples=90, smooth]
    plot (\x, {1 + 0.55*\x - 0.5*\x*\x + 0.28*\x*\x*\x});
  \node[blue!75!black, font=\scriptsize, anchor=south east] at (-0.85,1.05) {$f$};

  \fill (0,1) circle (1.1pt);

  \node[font=\scriptsize, align=center, anchor=north] at (0.3,-0.62)
    {each order matches one more derivative at $x_0$,\\
     and the error drops from $O(|x|)$ to $O(|x|^2)$ to $O(|x|^3)$};
\end{tikzpicture}
\end{center}

In several variables the same picture holds with \qt{tangent line} replaced by \emph{tangent
plane} and \qt{parabola} by a paraboloid, bowl-shaped or saddle-shaped according to the
signs of the eigenvalues of $\Hess f(x_0)$. That is exactly the classification carried out in
\cref{sec:hessian_test}.

% Source: Corsin Nick/Class Notes/Week 4.pdf, p. 12
% (Re-cited: the reading of the three terms and the figure above are additions.)
\begin{question}
For $f \in C^k(U, \mathbb{R})$, do the Taylor expansions of $t \mapsto f(x_0+th)$ in $t$ and of
$(th) \mapsto f(x_0+th)$ in $(th)$ around $0$ coincide?
\end{question}

\begin{answer}
\textbf{Yes.} This is, in fact, how we prove the multi-dimensional Taylor approximation, and why
we need the assumption $x_0 + th \in U$ for all $t \in [0,1]$.
\end{answer}

% Source: Corsin Nick/Class Notes/Week 4.pdf, p. 13
\begin{exercise}[Proof of second-order Taylor formula]
\label{ex:hessian_taylor_proof}
Prove the formula from above:
\[f(x_0+h) = f(x_0) + \nabla f(x_0)h + \tfrac{1}{2}h\transp \Hess f(x_0)h + \dots\]
for $f \in C^k(U,\mathbb{R})$ with $x_0 + th \in U$ for all $t \in [0,1]$.
\end{exercise}

% Kept inline: solved directly in Corsin Nick/Class Notes/Week 4.pdf, p. 13.
\begin{exercisesolution}[Proof of \cref{ex:hessian_taylor_proof}]
We expand $t \mapsto f(x_0+th)$ in $t$:
\[
\begin{aligned}
f(x_0+th) &= f(x_0) + t\left.\frac{d}{dt}\right|_{t=0}f(x_0+th) + \frac{t^2}{2}\left.\frac{d^2}{dt^2}\right|_{t=0}f(x_0+th) + o(|th|^2) \\
&= f(x_0) + t\sum_{i=1}^{n}\partial_i f(x_0)h_i + \frac{t^2}{2}\frac{d}{dt}\left.\sum_{i=1}^{n}\partial_i f(x_0+th)h_i\right|_{t=0} + o(|h|^2) \\
&= f(x_0) + t\nabla f(x_0)h + \frac{t^2}{2}\sum_{i,j=1}^{n}\partial_i\partial_j f(x_0)h_ih_j + o(|th|^2) \\
&= f(x_0) + t\nabla f(x_0)h + \frac{t^2}{2}h\transp \Hess f(x_0)h + o(|th|^2)
\end{aligned}
\]
Setting $t = 1$ gives the result.
\end{exercisesolution}


% --- exercises relocated here from the dissolved sheets ---

% Originally: content/week-05.tex, \exercisesheet{5}, problem 5.3
% Source: exercises/Ex5_Analysis2_eng.pdf

\begin{exercise}[Multiple choice (Landau notation) \textnormal{(important)}]
\label{ex:5.3}
Mark all and only the statements which are true.
\begin{enumerate}[label=\textbf{(\alph*)}]
  \item As $|x| \to 0$, if $f(x) = x_2x_1^2 + O(x_1^2)$, then $f(x) = O(|x|^2)$.
  \item As $x \to 0$, if $f(x) = x_2x_1^2 + O(x_1^3)$, then $f(x) = O(x_1^3)$.
  \item As $|x| \to 0$, if $f(x) = x_2x_1 + O(|x|^2)$, then $f$ is differentiable at $x = 0$.
  \item As $x \to 0$, if $f(x) = x_2x_1 + O(|x|^2)$, then $f$ is twice differentiable around
    $x = 0$.
  \item As $|x| \to 0$, if $f(x) = x_2x_1 + O(|x|^3)$, then $f$ is twice differentiable around
    $x = 0$ and $\partial_{11}f(0) = 0$.
\end{enumerate}
\exhint[Official hint]{Revise \textit{Corollary 10.45}.}
\exinfo{This exercise is Problem 5.3 of Problem Sheet 5, Analysis II, Spring Semester 2026.}
\end{exercise}

% Originally: content/exercise-sheets/week-05.tex, problem 5.5 --- never previously built
% Source: exercises/Ex5_Analysis2_eng.pdf

\begin{exercise}[Multiple choice (cost of a Taylor expansion) \textnormal{(optional)}]
\label{ex:5.5}
Assume $g \in C^\infty(\mathbb{R})$ and $f \in C^\infty(\mathbb{R}^n)$ is such that
$$f(x) = x_1 + x_2x_1 + O(|x|^4) \quad \text{as } x \to 0.$$
We want to compute the Taylor expansion of $(g\circ f)$ at $x = 0$ of the highest possible order
with the information we have. We can ask for the value $g^{(k)}(t)$ for arbitrary
$k \in \mathbb{N}$, $t \in \mathbb{R}$, but we have to pay 5 CHF each time.

For a general $g$, which is the degree of the best (i.e.\ highest-order) Taylor polynomial we can
compute? How expensive will it be to compute it?
\begin{enumerate}[label=\textbf{(\alph*)}]
  \item degree 3 and 20 CHF
  \item degree 2 and 15 CHF
  \item degree 2 and 10 CHF
  \item degree 3 and 15 CHF
\end{enumerate}
\exinfo{This exercise is Problem 5.5 of Problem Sheet 5, Analysis II, Spring Semester 2026.}
\end{exercise}

% Originally: content/exercise-sheets/week-05.tex, problem 5.6 --- never previously built
% Source: exercises/Ex5_Analysis2_eng.pdf

\begin{exercise}[Polynomials \textnormal{(optional --- ``interesting, but irrelevant for the course'')}]
\label{ex:5.6}
Let $P, Q \in \mathbb{R}[x_1,\dots,x_n]$ be polynomials of $n$ variables; assume that there are
positive numbers $M, \sigma$ such that
$$|P(x) - Q(x)| \leq M|x|^\sigma \quad \text{for all } |x| \leq 1.$$
Show that $P$ and $Q$ have the same coefficients of order smaller than $\sigma$. That is, if we
write $P(X) = \sum_{\alpha\in\mathbb{N}^n} a_\alpha X^\alpha$,
$Q(X) = \sum_{\alpha\in\mathbb{N}^n} b_\alpha X^\alpha$, then $a_\alpha = b_\alpha$ for all
$|\alpha| < \sigma$.
\exinfo{This exercise is Problem 5.6 of Problem Sheet 5, Analysis II, Spring Semester 2026.}
\end{exercise}

% Source: Jérôme Paschoud/Notizen/Woche 5.1 Taylor.pdf, p. 3
% Generator: Gemini 3.6 Flash (High)
\begin{exercise}[Evaluating partial derivatives using multi-index notation]
\label{ex:multi_index_evaluation}
Let $f(x,y,z) := x^3+2y+e^{2x}\cos(z)$. Calculate the partial derivative $\partial^\alpha f$ for the multi-index $\alpha = (2,0,1)$.
\end{exercise}

\begin{exercisesolution}[\cref{ex:multi_index_evaluation}]
The multi-index $\alpha = (2,0,1)$ corresponds to differentiating twice with respect to $x$, zero times with respect to $y$, and once with respect to $z$. Since $f$ is smooth, the order of differentiation does not matter.
\begin{align*}
\partial_x f &= 3x^2 + 2e^{2x}\cos(z) \\
\partial_{xx} f &= 6x + 4e^{2x}\cos(z) \\
\partial_{xxz} f &= -4e^{2x}\sin(z).
\end{align*}
Hence $\partial^{(2,0,1)} f(x,y,z) = -4e^{2x}\sin(z)$.
\end{exercisesolution}

% Source: Jérôme Paschoud/Notizen/Woche 5.1 Taylor.pdf, pp. 3--4
% Generator: Gemini 3.6 Flash (High)
\begin{exercise}[Taylor expansion of $\sin(xy)$]
\label{ex:taylor_expansion_sin_xy}
Compute the Taylor expansion of $f(x,y) := \sin(xy)$ to order $3$ around $(0,0)$. 
\end{exercise}

\begin{exercisesolution}[\cref{ex:taylor_expansion_sin_xy}]
We can solve this in two ways:
\begin{enumerate}[label=\textbf{(\alph*)}]
  \item \textbf{Using the general formula:}
    We list the multi-indices $\alpha$ with $|\alpha| \leq 3$ and compute the partial derivatives at $(0,0)$:
    \begin{itemize}
      \item $|\alpha|=0$: $f(0,0) = \sin(0) = 0$.
      \item $|\alpha|=1$: $\partial_x f = y\cos(xy) \implies 0$. $\partial_y f = x\cos(xy) \implies 0$.
      \item $|\alpha|=2$: $\partial_{xx} f = -y^2\sin(xy) \implies 0$. $\partial_{yy} f = -x^2\sin(xy) \implies 0$. $\partial_{xy} f = \cos(xy) - xy\sin(xy) \implies 1$.
      \item $|\alpha|=3$: $\partial_{xxx} f = -y^3\cos(xy) \implies 0$. $\partial_{yyy} f = -x^3\cos(xy) \implies 0$. $\partial_{xxy} f = -2y\sin(xy) - xy^2\cos(xy) \implies 0$. $\partial_{xyy} f = -2x\sin(xy) - x^2y\cos(xy) \implies 0$.
    \end{itemize}
    The only non-zero term is $\partial^{(1,1)} f(0,0) = 1$. The corresponding factor in the formula is $\frac{1}{\alpha!}h^\alpha = \frac{1}{(1!)(1!)}xy = xy$. Thus,
    \[f(x,y) = xy + o(|(x,y)|^3).\]
  \item \textbf{Using 1D expansions (the clever way):}
    We know that $\sin(t) = t - \frac{t^3}{6} + O(t^5)$ as $t \to 0$. Substituting $t = xy$ yields
    \[\sin(xy) = xy - \frac{(xy)^3}{6} + O((xy)^5) = xy + o(|(x,y)|^3).\]
\end{enumerate}
\end{exercisesolution}

% Source: Jérôme Paschoud/Notizen/Woche 5.1 Taylor.pdf, p. 4
% Generator: Gemini 3.6 Flash (High)
\begin{exercise}[Taylor expansion of $\sqrt{1+x+y^2}$]
\label{ex:taylor_expansion_sqrt_1_plus_x_plus_y2}
Compute the Taylor expansion of $f(x,y) := \sqrt{1+x+y^2}$ to order $2$ around $(0,0)$.
\end{exercise}

\begin{exercisesolution}[\cref{ex:taylor_expansion_sqrt_1_plus_x_plus_y2}]
Using the 1D expansion $\sqrt{1+t} = 1 + \frac{1}{2}t - \frac{1}{8}t^2 + O(t^3)$ with $t = x+y^2$, we obtain:
\begin{align*}
\sqrt{1+x+y^2} &= 1 + \frac{1}{2}(x+y^2) - \frac{1}{8}(x+y^2)^2 + O(|(x,y)|^3) \\
&= 1 + \frac{1}{2}x + \frac{1}{2}y^2 - \frac{1}{8}(x^2+2xy^2+y^4) + o(|(x,y)|^2) \\
&= 1 + \frac{1}{2}x - \frac{1}{8}x^2 + \frac{1}{2}y^2 + o(|(x,y)|^2).
\end{align*}
\end{exercisesolution}

% Supplement: Analysis_II_Script_v1.pdf, p. 59
\begin{exercise}[Taylor expansion of $\exp(\arctan(x-y))$]
\label{ex:taylor_expansion_exp_arctan}
Compute the Taylor expansion of $f(x,y) := \exp(\arctan(x-y))$ to order $2$ around $(0,0)$.
\end{exercise}

\begin{exercisesolution}[\cref{ex:taylor_expansion_exp_arctan}]
% Generator: Claude Sonnet 5 (High)
Write $s := x-y$, itself of order $|h|$. Since $\arctan$ is odd, $\arctan(t) = t + O(t^3)$, so
\[\arctan(s) = s + O(|h|^3) = s + o(|h|^2).\]
Since $\exp$ is smooth and bounded near $0$, an $o(|h|^2)$ correction inside $\exp$ contributes
only an $o(|h|^2)$ correction to the output. Consequently, to order $2$, $\exp(\arctan(s))$ and
$\exp(s)$ agree:
\[\exp(\arctan(x-y)) = \exp(x-y) + o(|(x,y)|^2).\]
Now $\exp(s) = 1+s+\tfrac12s^2+O(s^3)$ with $s^2 = x^2-2xy+y^2$, so
\[\exp(\arctan(x-y)) = 1 + x - y + \tfrac12x^2 - xy + \tfrac12y^2 + o(|(x,y)|^2).\]
\end{exercisesolution}

% Supplement: Analysis_II_Script_v1.pdf, p. 59
\begin{exercise}[Taylor expansion of $1/(1-x^2-y^2)$]
\label{ex:taylor_expansion_one_over_1_minus_x2_y2}
Compute the Taylor expansion of $f(x,y) := \dfrac{1}{1-x^2-y^2}$ to order $2$ around $(0,0)$.
\end{exercise}

\begin{exercisesolution}[\cref{ex:taylor_expansion_one_over_1_minus_x2_y2}]
% Generator: Claude Sonnet 5 (High)
Using $\tfrac{1}{1-t} = 1+t+O(t^2)$ with $t := x^2+y^2$, itself already $O(|h|^2)$, so
$t^2 = O(|h|^4) = o(|h|^2)$:
\[\frac{1}{1-x^2-y^2} = 1 + (x^2+y^2) + o(|(x,y)|^2).\]
\end{exercisesolution}

\begin{remark}
Together with \cref{ex:taylor_expansion_sin_xy,ex:taylor_expansion_sqrt_1_plus_x_plus_y2} above,
these two exercises are exactly script Exercise 10.36: \qt{compute the Taylor polynomials up
to the quadratic order of $\sin(xy)$, $\sqrt{1+x+y^2}$, $\exp(\arctan(x-y))$,
$1/(1-x^2-y^2)$ --- don't use the general formula!} All four are worked by the same \qt{1D
substitution} technique this section uses throughout.
\end{remark}

% Source: Jérôme Paschoud/Notizen/Woche 5.1 Taylor.pdf, p. 4
% Generator: Gemini 3.6 Flash (High)
\begin{exercise}[Taylor expansion of $\frac{\cos x}{1-y^2}$]
\label{ex:taylor_expansion_cos_over_1_minus_y2}
Compute the Taylor expansion of $f(x,y) := \frac{\cos x}{1-y^2}$ to order $4$ around $(0,0)$.
\end{exercise}

\begin{exercisesolution}[\cref{ex:taylor_expansion_cos_over_1_minus_y2}]
We multiply the 1D expansions for $\cos x$ and $\frac{1}{1-y^2}$.
As $x \to 0$, $\cos x = 1 - \frac{x^2}{2} + \frac{x^4}{24} + O(x^6)$.
As $y \to 0$, $\frac{1}{1-y^2} = 1 + y^2 + y^4 + O(y^6)$.
Multiplying these together and keeping terms up to order $4$:
\begin{align*}
\frac{\cos x}{1-y^2} &= \left(1 - \frac{x^2}{2} + \frac{x^4}{24} + O(x^6)\right)\left(1 + y^2 + y^4 + O(y^6)\right) \\
&= 1 + y^2 + y^4 - \frac{x^2}{2} - \frac{x^2y^2}{2} + \frac{x^4}{24} + o(|(x,y)|^4) \\
&= 1 - \frac{1}{2}x^2 + y^2 + \frac{1}{24}x^4 - \frac{1}{2}x^2y^2 + y^4 + o(|(x,y)|^4).
\end{align*}
\end{exercisesolution}

% Source: old_exams/fs2023/Prüfung.pdf (9 August 2021), Teil B, Aufgabe 1, p. 13
% Extractor: Gemini 3.6 Flash (High)
% Correction: Claude Opus 5 (High) --- part (c) restored to the source's "all critical points";
% Flash had weakened it to "in a neighbourhood of x = 0", and then answered neither version. Part
% (d) restored: the source asks about x = 1 as well, and that is the interesting one.
\begin{exercise}[Taylor expansion of a parametric integral]
\label{ex:taylor_parametric_integral}
Consider the function $f : \mathbb{R} \to \mathbb{R}$ defined by the parametric integral
\[f(x) := \int_{x-1}^{x+1} \cos(\pi t^2) \, dt.\]
The integral need not be evaluated explicitly.
\begin{enumerate}[label=\textbf{(\alph*)}]
  \item Compute $f'(x)$, and simplify it to a product of two sines.
  \item Find the coefficients of the fourth-order Taylor polynomial of $f$ around $x = 0$,
  \[f(x) = f(0) + a_1 x + a_2 x^2 + a_3 x^3 + a_4 x^4 + O(x^5).\]
  \item Determine \emph{all} critical points of $f$.
  \item Decide, for each of the critical points $x = 0$ and $x = 1$, whether it is a local maximum,
    a local minimum, or neither.
\end{enumerate}
\exinfo{This exercise is taken from the exam of 9 August 2021, Part B, Problem 1. The original
opens with a part asking whether $f$ is continuously differentiable, bounded, and even.}
\end{exercise}

